% ============================================================
% Formulario, argomento 6: pmf condizionate e funzioni di una variabile (dispensa, capitolo 6)
% ============================================================
\argomento{Pmf condizionate e funzioni di una variabile}{Dispensa: capitolo 6}[Pmf condizionate e funzioni]

\voce{pmfcond}{Pmf condizionata a un evento}{$A$ evento con $P(A) \neq 0$}
  {p_{X \mid A}(x) = P(X = x \mid A) = \frac{P(\{X = x\} \cap A)}{P(A)}}
  {\variante{se l'evento è definito da $X$ stessa, $A = \{X \in B\}$}
            {p_{X \mid X \in B}(x) = \frac{p_X(x)}{P(X \in B)} \ \text{per } x \in B, \quad 0 \ \text{altrove}}
   \variante{è una pmf}{p_{X \mid A}(x) \geq 0, \quad \sum_{x \in \mathcal{X}} p_{X \mid A}(x) = 1}}

\voce{totalepmf}{Probabilità totale per le pmf}{$\{E_i\}$ partizione di $\Omega$, $P(E_i) > 0$}
  {p_X(x) = \sum_{i=1}^{M} p_{X \mid E_i}(x)\, P(E_i)}
  {\variante{al contrario, Bayes: il caso sapendo il valore}
            {P(E_i \mid X = x) = \frac{p_{X \mid E_i}(x)\, P(E_i)}{p_X(x)}}}

\voce{mediacond}{Media condizionata a un evento}{$P(E) > 0$}
  {E[X \mid E] = \sum_{x \in \mathcal{X}} x\, p_{X \mid E}(x)}
  {}

\voce{mediamedie}{Media come media delle medie condizionate}{$\{E_m\}$ partizione di $\Omega$}
  {E[X] = \sum_{m=1}^{M} P(E_m)\, E[X \mid E_m]}
  {\variante{due casi, $E$ e $\compl{E}$}{E[X] = P(E)\, E[X \mid E] + P(\compl{E})\, E[X \mid \compl{E}]}}

\voce{pmffunzione}{Pmf di una funzione $Y = g(X)$}{alfabeto di $Y$: $\mathcal{Y} = g(\mathcal{X})$}
  {p_Y(y) = \sum_{x \,:\, g(x) = y} p_X(x)}
  {\variante{$g$ biunivoca (solo una ridenominazione)}{p_Y(y) = p_X\!\left(g^{-1}(y)\right)}}

\voce{valoremedio}{Teorema fondamentale per la media}{teorema del valore medio}
  {E[Y] = E[g(X)] = \sum_{x \in \mathcal{X}} g(x)\, p_X(x)}
  {\variante{per esempio}{E[X^2] = \sum_{x \in \mathcal{X}} x^2\, p_X(x)}
   \variante{in generale}{E[g(X)] \neq g(E[X])}}
